\documentclass[11pt]{article}
\usepackage[margin=1.05in]{geometry}
\usepackage{amsmath,amsthm,amssymb,mathtools,bm}
\usepackage{booktabs,array,enumitem}\usepackage[expansion=false]{microtype}
\usepackage{xcolor}
\usepackage[colorlinks=true,linkcolor=blue!60!black,citecolor=blue!60!black,urlcolor=blue!60!black]{hyperref}
\theoremstyle{plain}
\newtheorem{theorem}{Theorem}[section]\newtheorem{proposition}[theorem]{Proposition}\newtheorem{lemma}[theorem]{Lemma}
\newtheorem{corollary}[theorem]{Corollary}\newtheorem{conjecture}[theorem]{Conjecture}
\theoremstyle{definition}\newtheorem{definition}[theorem]{Definition}\newtheorem{remark}[theorem]{Remark}
\newtheorem{example}[theorem]{Example}
\newcommand{\R}{\mathbb R}\newcommand{\Z}{\mathbb Z}\newcommand{\E}{\mathbb E}\newcommand{\Var}{\operatorname{Var}}\newcommand{\Cov}{\operatorname{Cov}}
\newcommand{\relu}{\operatorname{relu}}\newcommand{\diag}{\operatorname{diag}}\newcommand{\tr}{\operatorname{tr}}
\newcommand{\Tr}{\operatorname{Tr}}\newcommand{\1}{\mathbf 1}\newcommand{\cA}{\mathcal A}\newcommand{\cB}{\mathcal B}
\newcommand{\cF}{\mathcal F}\newcommand{\cT}{\mathcal T}\newcommand{\cH}{\mathcal H}\newcommand{\aff}{\mathfrak{aff}}
\title{\textsc{Localization in the commutant, stage 11:}\\ \large tilings, Bratteli diagrams and singular foliations.\\
Depth as a parabolic renormalization, the two tiers as one flow,\\ and the walls of a \textsc{ReLU} network as a stratified foliation}
\author{}
\date{10 October 2026}
\begin{document}
\maketitle
\begin{abstract}
We read eight papers: the tiling-groupoid and Bratteli-diagram papers of Bellissard--Julien--Savinien and Julien--Savinien, the
bi-infinite Bratteli diagrams and translation surfaces of Putnam--Trevi\~no, the measure-free Koopman--von Neumann spectral triples
of Giannakis--Montgomery, Paul's semiclassics beyond Ehrenfest time, and three papers on singular foliations (Francis; Fischer and
Laurent-Gengoux; Louis). We use them to place the bias-free He-initialised \textsc{ReLU} network with Gaussian input, and the
two-tier (localization over network) framework of stages 9 and 10, inside tiling and foliation theory. The results are these.
(i)~At a fixed input the network is a weighted Bratteli diagram on the layered neuron set. Forward activations and backward
sensitivities are its two harmonic states, and Putnam--Trevi\~no's conservation of $\sum_v\nu_r\nu_s$ is backpropagation: the
output equals $\langle h_l,\partial_{h_l}f\rangle$ at every layer. The \textsc{ReLU} rescaling symmetry is the diagonal
(Teichm\"uller-type) flow, and Rauzy--Veech induction is itself a gated linear network whose activation tiling is the partition into
Rauzy cylinders. (ii)~Depth is a \emph{parabolic} renormalization. The folding map of the arc-cosine kernel is the germ
$\theta-\theta^2/3\pi-\theta^3/18\pi^2+\dots$, with Fatou coordinate $3\pi/\theta+\tfrac32\log\theta$ and iterative residue $3/2$. Hence
$\theta_l=3\pi/(l+\tfrac32\log l+4.39+o(1))$ from orthogonal inputs. At $l=16$ this gives $0.384$ against the exact $0.378$, whereas
the stage-10 law $3\pi/l$ gives $0.589$. (iii)~The two tiers are one flow. The unresolved gate variance at depth $l$ and localization
time $t$ depends only on the Fatou time $\tau=\Phi(\theta_0(t))+l-1$. Localization deforms the state, a layer shifts the diagram,
and the Fatou coordinate is the renormalization time, exactly as in Putnam--Trevi\~no's description of the Teichm\"uller flow. One
layer is worth $\sqrt2/3\pi$ of $\sqrt{1+t}$, and measured width-512 data collapse on $\tau$ except at the last layer. (iv)~The walls form
a stratified singular foliation whose leaves are the faces of the tiling. Crossings of walls of one layer are normal crossings with
abelian isotropy. A deep wall crossing a wall that feeds it is bent; its formal transverse model is a three-line arrangement whose
isotropy algebra is the non-abelian $\aff(1)$, and the bending is the gated transport coefficient at which Price's theorem gives
birth to the third cumulant. Faces are contractible and every wall has a global defining function, so outer holonomy is trivial and
the noncommutativity is all in this isotropy. (v)~The upper tier is a semiclassical deformation: averaging over Gaussian balls is a
coherent-state (Husimi) symbol with an exact $\sqrt\hbar$ boundary layer at every wall. At the target scale there are $n/4$ unresolved
gates per layer, so the infinite-width description is a kinetic (van Hove) limit rather than a semiclassical one. (vi)~The mean is a
border functional on the codimension-one faces, written as a supertrace pairing of the wall flip with the normal derivative.
Every claim is labelled proved, known, conjectural or an analogy. The identities are checked by \texttt{scripts/verify\_stage11.py}.
\end{abstract}
\tableofcontents
\section{The placement in one page}
\label{sec:summary}
Notation as in stage 10: $W_l\in\R^{n\times n}$ with i.i.d.\ $N(0,2/n)$ entries, $h_0=x\sim\gamma=N(0,I_n)$, $z_l=W_lh_{l-1}$,
$h_l=\relu(z_l)$, gates $g_l=\1\{z_l>0\}$, $D_l=\diag(g_l)$; for an output functional $f=c^{\top}h_L$ put $\delta_l=\partial f/\partial h_l$.
The activation tiling at depth $l$ is the partition of input space by $(g_1,\dots,g_l)$; its walls are $Z_{l,a}=\{z_{l,a}=0\}$. The
upper tier is the family of laws $N(m,\hbar I)$ (or Eldan's path with $\hbar=1/(1+t)$), and $F(m,\hbar)=\E f(m+\sqrt\hbar Z)$.

\paragraph{1. The cocycle (Section~\ref{sec:cocycle}).} At a fixed input the network is a weighted Bratteli diagram on the layered
neuron set, with edge matrices $E_l(x)=D_l(x)W_l$. Forward activations are a left-harmonic state, backward sensitivities a right-harmonic
state, and Putnam--Trevi\~no's invariant $\sum_{v\in V_l}\nu_r(v)\nu_s(v)$ is $\langle h_l,\delta_l\rangle=f(x)$ at every layer
(Theorem~\ref{thm:bratteli}). Neuron $i$ of layer $l$ is a rectangle of width $h_{l,i}$ and height $\delta_{l,i}$ (its gradient$\times$input
attribution), each layer is a zippered-rectangles presentation of the same area, and the \textsc{ReLU} rescaling symmetry is the diagonal
flow that fixes every rectangle's area. Moving the input across a wall is a rank-one elementary move of the cocycle, the analogue of a
Rauzy move. Conversely, Rauzy--Veech induction is a gated linear network: one gate per step, its activation regions are the Rauzy
cylinders, and its Jacobians are the Rauzy--Veech cocycle (Proposition~\ref{prop:rv}).

\paragraph{2. Depth is parabolic (Section~\ref{sec:parabolic}).} The pair angle obeys $\theta_{l}=F(\theta_{l-1})$ with
$F(\theta)=\theta-\theta^2/3\pi-\theta^3/18\pi^2+O(\theta^4)$, a parabolic germ. Its Fatou coordinate is
$\Phi(\theta)=3\pi/\theta+\tfrac32\log\theta+O(\theta)$ (iterative residue $3/2$), and $\Phi\circ F=\Phi+1$. This replaces the stage-10
folding law $3\pi/l$, which is $56\%$ off at $l=16$, by $\theta_l=3\pi/(l+\tfrac32\log l+4.39+o(1))$, which is $1.6\%$ off
(Theorem~\ref{thm:fatou}). The He-critical network is a Farey-type (additive, parabolic) renormalization, not a Gauss-type (hyperbolic)
one. The infinite sum $\sum_l\theta_l$ diverges, so at infinite width two distinct inputs disagree on gates at infinitely many layers.

\paragraph{3. The two tiers are one flow (Theorem~\ref{thm:collapse}).} The unresolved gate variance at depth $l$ and localization time
$t$ is a function of the Fatou time $\tau=\Phi(\theta_0(t))+l-1$ alone. In Putnam--Trevi\~no's words, the Teichm\"uller flow is a
deformation of the state followed by a shift of the diagram: here the localization deforms the state, a layer shifts the diagram, and
the Fatou coordinate measures the renormalization time. One layer of depth is worth $\Delta\sqrt{1+t}=\sqrt2/3\pi\approx0.150$. The
measured width-512 disagreement fractions of stage 10, sorted by $\tau$, fall on one curve except at the last layer.

\paragraph{4. The walls are a stratified singular foliation (Section~\ref{sec:walls}).} Its leaves are the open faces of the tiling. Walls
of one layer cross normally, so the isotropy there is abelian ($b$-calculus, Debord, Helffer--Nourrigat cone equal to the log
cotangent). A deep wall crossing a wall that feeds it is bent by the gated transport coefficient $T=\partial z_{l,a}/\partial h_{k,b}$.
Its formal transverse model is the three-line arrangement $y_1y_2(y_2+Ty_1)=0$, whose isotropy algebra is $\aff(1)$ with
$[E,\vartheta]=\vartheta$. The linear isotropy drops from dimension 2 to 1 (Theorem~\ref{thm:bent}). The bent crossings are where Price's theorem gives birth to the third cumulant: with
orthogonal first-layer rows, $\kappa_3(z_{2,a})=c_3\sum_bT_{ab}^3|w_b|^3$, $c_3=(\pi+2)/(2\pi)^{3/2}$. Faces are contractible and walls
have global defining functions, so by Fischer--Laurent-Gengoux and Francis there is no holonomy. The wall foliation's
noncommutativity is purely isotropic and sits exactly on the memory.

\paragraph{5. The upper tier is a semiclassical deformation (Section~\ref{sec:deformation}).} $N(m,\hbar I)$ is a coherent state and
$F(m,\hbar)$ the network's Husimi symbol. Near a wall, $F(m,\hbar)-f(m)=\sqrt\hbar\,\kappa\,\psi(d/\sqrt\hbar)$ up to exponentially small
terms, with $\psi(u)=\varphi(u)-|u|\Phi_{\mathcal N}(-|u|)$. So $F$ is smooth on the parabolic blowup of the walls, and corners need
iterated blowups along the face poset (Theorem~\ref{thm:layer}). The target $\hbar=1$ is deep in the quantum regime ($n/4$ unresolved gates
per layer). The infinite-width (moment-chain) description is a weak-coupling, high-density kinetic limit. In Giannakis--Montgomery's
language, chains are Fock truncations, deep layers approach the commutative (classical) quotient, and the target is a kernel mean
embedding whose estimation error is a maximum mean discrepancy.

\paragraph{6. Borders (Section~\ref{sec:borders}).} The depth tower carries two Bratteli diagrams. The tile diagram refines and the fiber
diagram folds; the fibers of $h_l$ play the role of Bellissard--Julien--Savinien's supertiles. Border dimension is the codimension of the
face containing a point. The mean is a border functional, $\E f=(2\pi)^{-1/2}\sum_{\text{codim-1 faces}}\kappa_F\alpha_F$, equivalently the
supertrace pairing of the wall flip with the normal derivative (checked to eight digits in two dimensions).

\section{What the eight papers do}
\label{sec:readings}
\paragraph{Bellissard--Julien--Savinien \cite{BJS}.} A tiling of finite local complexity has a hull $\Omega$ (the closure of its
translates) and a canonical transversal $\Xi$ (tilings with a puncture at the origin), a Cantor set carrying the étale equivalence relation
$R_\Xi$ of translation. A \emph{proper nested sequence} of tilings $T_0\prec T_1\prec\dots$ (each tile of $T_{n+1}$ a patch of $T_n$, with
buffer zones) gives a Bratteli diagram whose vertices at level $n$ are the collared prototiles of $T_n$ and whose edges record which tile
sits in which supertile. Infinite paths are points of $\Xi$ (the Robinson map), and tail equivalence of paths is an AF subrelation
$R_{AF}\subsetneq R_\Xi$. The two agree on a dense $G_\delta$ of full measure (points that are eventually deep inside supertiles) and
differ on a dense, null $F_\sigma$: tilings whose origin stays on supertile boundaries at every scale. To recover $R_\Xi$ the authors add
\emph{horizontal} edges (adjacency of supertiles through collars) labelled by translation vectors. In dimension one this reduces to
$R_\Xi=R_{AF}\vee\{(x^i_{\min},x^i_{\max})\}$, the Vershik pairs of maximal and minimal paths (Fibonacci: two pairs; Thue--Morse: four).
\emph{Essence:} the coarse-graining tower is a Bratteli diagram, generic points see only the AF part, and the full dynamics is the AF part
plus a thin set of boundary identifications, which carry all the noncommutativity that the AF part misses.

\paragraph{Julien--Savinien \cite{JS}.} For substitution tilings the boundary identifications are made combinatorial. The substitution
induces substitutions $\omega_j$ on decorated $j$-faces for every $j<d$, and each has its own Bratteli diagram $B^j$. A multi-diagram joins
them by horizontal edges (a $j$-face lies on the boundary of a $(j{+}1)$-face) and by \emph{escaping} edges (a face vanishes into the interior
of a supertile). A path in $B^d$ has \emph{borders}: tails in $B^j$ obtained by following horizontal edges. Its border dimension is the least
such $j$, geometrically the dimension of $\bigcap_n\lambda^n\omega^{-n}(T)^j$. Two paths are equivalent iff they have the same border
dimension $j$ and tail-equivalent $j$-borders, giving $R_\Xi=\bigsqcup_jR^j_B$ with $R^d_B=R_{AF}$. Each $R^j_B$ with $j<d$ lives on a dense
null set and is not AF. \emph{Essence:} the full relation decomposes by the dimension of the face on which a point is trapped across all
scales, and each stratum is a tail relation of the face substitution of that dimension.

\paragraph{Putnam--Trevi\~no \cite{PT}.} A bi-infinite Bratteli diagram has levels indexed by $\Z$, two partial orders $\le_s$ (edges
with a common source) and $\le_r$ (common range), and a state $(\nu_r,\nu_s)$. Here $\nu_r(v)=\sum_{r(e)=v}\nu_r(s(e))$ propagates forward,
$\nu_s(v)=\sum_{s(e)=v}\nu_s(r(e))$ propagates backward, and $\sum_{v\in V_n}\nu_r(v)\nu_s(v)$ is independent of $n$ (their Lemma 2.7). The
cylinder of paths through a finite path $p$ has measure $\nu_r(s(p))\nu_s(r(p))$ and maps to a rectangle of width $\nu_r(s(p))$ and
height $\nu_s(r(p))$. After removing extremal paths and the \emph{singular} set $\Sigma_B=\{\Delta_s\Delta_r\ne\Delta_r\Delta_s\}$, where the
two successor maps fail to commute, the quotient of the path space by both successor relations is a translation surface. Right-tail
equivalence becomes the horizontal foliation and left-tail equivalence the vertical one. Their $C^*$-algebras are AF and foliation
algebras linked by explicit subhomogeneous inductive limits. The successor flip $F_B$ gives an even Fredholm module whose commutant
inside the AF algebra is the algebra of the quotient (Theorem 9.7), and an excision argument computes the $K$-theory. Infinite
genus appears exactly when $\Sigma_B\neq\emptyset$ (Chamanara's surface). In finite genus the diagram comes from backward Rauzy--Veech
induction, $K_0(C^*(F^+))\cong H_1(S,\Sigma;\Z)$, the order on $K_0$ is given by the Schwartzman asymptotic cycle (the top Zorich
direction), and the Teichm\"uller flow is a deformation of the state $(e^{-t}\nu_r,e^{t}\nu_s)$ followed by a shift of the diagram.
\emph{Essence:} two transverse tail structures and a pairing-conserving state make a flat surface. The obstruction to flatness is the
non-commutation of the two successor maps, and renormalization trades state deformation against shifts.

\paragraph{Giannakis--Montgomery \cite{GM}.} Koopman evolution $f\mapsto f\circ\Phi^t$ is studied without a reference measure, on a
reproducing-kernel Hilbert algebra $H$. The generator $V=S+A$ (symmetric plus antisymmetric) is dilated to a skew-adjoint tridiagonal
$L$ on $\ell^2(\mathbb N)\otimes H$ with $A$ on the diagonal and $\pm S$ off it. The Fibonacci vector $E f=\sum_iF_ie_i\otimes f$ intertwines,
$LE=EV$ (because $F_{i+1}-F_{i-1}=F_i$), so the corner $(p_1\otimes1)L^kE_n$ reproduces $V^k$. A second quantization on a weighted
symmetric Fock space (a Banach algebra under $\vee$ for subconvolutive weights) turns $L$ into a free derivation $D$. Creation and
annihilation operators are bounded, $(\mathcal A_n,\mathcal F,-iD)$ is a weak spectral triple, and $\Psi([D,a])=V(\Psi a)$. Points of $X$
embed as high-grade limits of truncated coherent states. The creation/annihilation commutators vanish on high grades, so the quotient
by the compact diagonal ideal is commutative (the classical limit), and expectations under any Borel measure are kernel mean embeddings
$\Psi_\mu=\int\Psi_x\,d\mu$. \emph{Essence:} a nonlinear or dissipative dynamics becomes linear and unitary on an amplified Fock space, the
classical system is the high-occupation quotient, and averages over laws are mean embeddings.

\paragraph{Paul \cite{Paul}.} For an Anosov Hamiltonian flow, the quantum evolution of an observable up to times $\hbar^{-2+\epsilon}$
(far beyond the Ehrenfest time $\log(1/\hbar)/\lambda$) is still described by a symbol. That symbol is an element of the crossed product of
$C(h^{-1}(I))$ by the group generating the strong unstable foliation, and it is obtained by pushing the initial symbol forward along
the flow. Operators are written as superpositions of rank-one maps between Lagrangian states built from Gaussian coherent states
along unstable leaves. Products of symbols are convolutions along leaves, and off-diagonal Husimi functions $R_\psi(z,z')$ with $z'$ on
the unstable leaf of $z$ are the matching states. \emph{Essence:} past the time at which a coherent state is sheared across many
scales, the classical limit is the noncommutative algebra of the foliation that the shearing produces, and two-point objects take over
from densities.

\paragraph{Francis \cite{Francis}.} A transverse order $k$ foliation has finitely many leaves: one hypersurface $L$ and the components
of its complement. It is generated near $L$ by vector fields tangent to $L$ to order $k$, locally $\{y^k\partial_y,\partial_{x_i}\}$. Loops
in $L$ have holonomy only at the level of $(k-1)$-jets, a homomorphism $\pi_1(L)\to J^{k-1}$, and this is a complete local invariant that
takes all values. The holonomy groupoid is a blowup of $(M\setminus L)^2\cup L^2$ in which $L^2$ is replaced by the gauge groupoid of a
principal $J^k_{\R}$-bundle. Its $C^*$-algebra is an extension $0\to\mathcal K(\oplus\mathcal K)\to C^*(G)\to C^*(\Gamma_{\R})\otimes\mathcal
K\to0$, and $J^2\cong$ the $ax+b$ group. If $L$ has a global defining function with a given $(k-1)$-jet (Scott's $b^k$-manifolds), the holonomy
is trivial. \emph{Essence:} singular foliations with the same leaves are distinguished by order of tangency and jet holonomy, and the
isotropy groups are solvable jet groups.

\paragraph{Fischer--Laurent-Gengoux \cite{FLG}.} Near a leaf $L$ with transverse model $\cT_0$ (a formal singular foliation of
$\R^d$ at $0$), singular foliations are classified formally by an outer holonomy $\pi_1(L)\to\mathrm{Out}(\cT_0)$ together with a
finite-dimensional principal bundle, built from a Galois cover and a group extension of $\mathrm{Inner}(\cT_0)$. If $L$ is contractible the
foliation is the product with $\cT_0$; if $L$ is simply connected only the principal bundle remains; if $\cT_0$ vanishes quadratically
only the outer holonomy remains. Parallel transport along a leaf maps transverse leaves to transverse leaves (Yang--Mills
connections). \emph{Essence:} ``travelling along a leaf one sees the same landscape through the window''. What can change is only the
global twisting, measured by symmetries of the transverse model.

\paragraph{Louis \cite{Louis}.} For a singular foliation $\cF$ with an anchored bundle $\rho:A\to TM$, the Nash blowup replaces each point by
all limits of the kernels $\ker\rho_x$ from regular points. The pulled-back foliation becomes Debord, i.e.\ the image of a Lie algebroid
$D_\cF$ (Mohsen's groupoid). The Helffer--Nourrigat cone $\mathrm{HN}(\cF)=\overline{\bigcup_{x\ \mathrm{regular}}\rho^*_x(T^*_xM)}$ is the image
of $D^*_\cF$; it is a union of symplectic leaves of the linear Poisson structure, and it is exactly where the principal symbol of a
longitudinal operator is intrinsic. Longitudinal ellipticity means positivity of the symbol on $D_\cF^*$. For Debord foliations the
blowup is trivial and $\mathrm{HN}=A^*$. \emph{Essence:} the symbol space of a singular foliation is the closure of its regular conormal
data, and resolving the foliation by limits of tangent spaces makes it a Lie algebroid.

\section{The cocycle: a network is a weighted Bratteli diagram}
\label{sec:cocycle}
\begin{definition}[The layered diagram of an input]
For $x$ off the walls let $B_x$ be the bi-infinite-in-spirit (here finite, levels $0,\dots,L$) weighted Bratteli diagram with
$V_l=[n]$ (neurons of layer $l$; $V_0$ the input coordinates) and edge weights $E_l(x)=D_l(x)W_l$ from $V_{l-1}$ to $V_l$. A path
$p=(p_0,\dots,p_L)$ has weight $\omega(p)=\prod_{k=1}^LE_k(x)_{p_kp_{k-1}}$. Put $\nu_r(l)=h_l$ and $\nu_s(l)=\delta_l$ for $f=c^\top h_L$.
\end{definition}
\begin{theorem}[Forward and backward states; the conserved pairing]\label{thm:bratteli}
On the tile of $x$:
\begin{enumerate}[leftmargin=*]
\item $\nu_r(l)=E_l\nu_r(l-1)$ and $\nu_s(l-1)=E_l^{\top}\nu_s(l)$: activations are a left-harmonic state and sensitivities a right-harmonic
state of $B_x$.
\item $\sum_{v\in V_l}\nu_r(v)\nu_s(v)=\langle h_l,\delta_l\rangle=f(x)$ for every $l$ (Putnam--Trevi\~no's Lemma 2.7).
\item $f(x)=\sum_{p}x_{p_0}\,\omega(p)\,c_{p_L}$, and the paths through a fixed sub-path $q$ from level $l$ to level $m$ contribute
$h_{l,q_l}\,\omega(q)\,\delta_{m,q_m}$: the cylinder measure $\nu_r(s(q))\,\omega(q)\,\nu_s(r(q))$.
\item For $s\in\R^n_{>0}$ the substitution $W_l\mapsto\diag(s)W_l$, $W_{l+1}\mapsto W_{l+1}\diag(s)^{-1}$ leaves all gates and $f$ unchanged and acts by
$h_l\mapsto s\odot h_l$, $\delta_l\mapsto s^{-1}\odot\delta_l$. It preserves every rectangle $R_{l,i}=[0,h_{l,i}]\times[0,\delta_{l,i}]$: the
\textsc{ReLU} rescaling symmetry is the diagonal (Teichm\"uller-type) flow $(e^{-u}\nu_r,e^{u}\nu_s)$, neuron by neuron.
\item Crossing the wall $Z_{l,a}$ changes $E_l$ by $\pm e_ae_a^{\top}W_l$ (rank one); $f$ is continuous and $\nabla f$ jumps by
$\delta_{l,a}\nabla z_{l,a}$, the kink of stage 10.
\end{enumerate}
\end{theorem}
\begin{proof}
(1) is the forward recursion and the chain rule $\delta_{l-1}=W_l^{\top}D_l\delta_l$. (2):
\[
\langle h_l,\delta_l\rangle=\langle D_lW_lh_{l-1},\delta_l\rangle=\langle h_{l-1},W_l^{\top}D_l\delta_l\rangle=\langle h_{l-1},\delta_{l-1}\rangle,\qquad\langle h_L,\delta_L\rangle=f.
\] (3) expands the matrix product. (4): $\relu(s_iz)=s_i\relu(z)$ for $s_i>0$, so
$z_{l+1}$ and all later quantities are unchanged; the chain rule gives the action on $\delta_l$. (5): only $g_{l,a}$ changes, $z_{l,a}=0$ on the
wall so $h_l$ is continuous, and later gates are unchanged at a generic wall point.
\end{proof}
Each layer is therefore a presentation of the same total area $f(x)$ by $n$ rectangles, re-cut by the cocycle $E_{l+1}$ in passing to the
next layer: widths transform by $E_{l+1}$ and heights by $E_{l+1}^{\top}$. This is the algebra of zippered rectangles under Rauzy--Veech
induction, where $\lambda'=\Theta^{-\top}\lambda$, $h'=\Theta h$ and $\lambda'\cdot h'=\lambda\cdot h$ \cite[\S12.2]{PT}. The area of rectangle
$R_{l,i}$ is the gradient$\times$input attribution of neuron $(l,i)$. Two differences are essential. The weights have both signs, so
the cylinder ``measure'' is signed and the quotient of the path space is not a flat surface. For a nonnegative network ($W\ge0$, e.g.\
the majorant $|W|$) Putnam--Trevi\~no's construction applies verbatim once orders on edges are fixed. And $E_l$ is neither unimodular nor
invertible, so the cocycle is critical rather than hyperbolic (Section~\ref{sec:parabolic}).

\begin{proposition}[Rauzy--Veech induction is a gated linear network]\label{prop:rv}
Let $\pi$ be an irreducible permutation on an alphabet $\mathcal A$ and let $R$ be one step of Rauzy--Veech induction on
$\lambda\in\R^{\mathcal A}_+$. The step is $\lambda\mapsto\lambda-g\,\lambda_{\alpha(1)}e_{\alpha(0)}-(1-g)\lambda_{\alpha(0)}e_{\alpha(1)}$ with the gate
$g=\1\{\lambda_{\alpha(0)}>\lambda_{\alpha(1)}\}$, the sign of a linear form, and the next permutation (the wiring) depends on $g$. Hence
$\lambda\mapsto R^k\lambda$ is piecewise linear and its linear pieces are the Rauzy cylinders, the sets of $\lambda$ with a given type
sequence $(g_1,\dots,g_k)$. On each cylinder the map is the inverse transpose of a product of $k$ elementary unimodular matrices (the
Rauzy--Veech cocycle). Each later gate is the sign of a linear form pulled back through the earlier steps, a bent hyperplane. So the
cylinder partition is the activation tiling of a depth-$k$ gated network with one gate per layer and gate-dependent wiring. Veech's
theorem \cite{Veech} (cylinders shrink to rays for Lebesgue-almost every $\lambda$) says that this network's gate sequence determines its
input projectively.
\end{proposition}
\begin{proof}
Direct from the definition of the step \cite[\S12.2]{PT}; the induction on $k$ is the chain rule for piecewise-linear maps.
\end{proof}
The check (Verification, item 2) runs $12$ steps for $d=4$, $\pi=(4321)$. The Jacobian on the cylinder is an integer matrix of
determinant $1$, perturbations staying in the cylinder obey the linear map to $2\times10^{-11}$, and all $2^6=64$ type patterns of length
$6$ occur. The comparison is instructive:
\begin{center}\small
\begin{tabular}{>{\raggedright\arraybackslash}p{4.4cm}>{\raggedright\arraybackslash}p{5.2cm}>{\raggedright\arraybackslash}p{5.2cm}}\toprule
 & Rauzy--Veech induction & He-\textsc{ReLU} network\\\midrule
gate & $\1\{\lambda_{\alpha(0)}>\lambda_{\alpha(1)}\}$, one per step & $\1\{z_{l,a}>0\}$, $n$ per layer\\
cocycle & elementary, unimodular, $\ge0$ & $D_lW_l$, signed, rank $\approx n/2$\\
contraction of cylinders & exponential after Zorich acceleration & parabolic: chords $\sim3\pi/(l+\tfrac32\log l)$\\
spectrum & Forni, Avila--Viana: $g$ positive exponents & critical: top exponent $\approx0$\\
order on $K_0$ & Schwartzman cycle (top Zorich direction) & collective direction of the layer covariance\\
invariant measure & Veech's Gauss measure & law of $h_l$ under $\gamma$\\\bottomrule
\end{tabular}
\end{center}

\section{Depth is a parabolic renormalization, and the two tiers are one flow}
\label{sec:parabolic}
Let $J(\theta)=\sin\theta+(\pi-\theta)\cos\theta$ and $F(\theta)=\arccos(J(\theta)/\pi)$ on $[0,\pi]$: the angle map of the arc-cosine kernel
\cite{CS}, under which the angle between $h_{l-1}(x)$ and $h_{l-1}(x')$ evolves at infinite width (stage 10, Theorem~\emph{folding}).
\begin{theorem}[Depth is a parabolic germ; Fatou coordinate]\label{thm:fatou}
\begin{enumerate}[leftmargin=*]
\item $F$ is real-analytic on $[0,\pi)$, it satisfies $F(\theta)<\theta$ on $(0,\pi]$, and
\[
F(\theta)=\theta-\frac{\theta^2}{3\pi}-\frac{\theta^3}{18\pi^2}+O(\theta^4)
\]
(numerically the $\theta^4$ coefficient is $-0.00767$). Thus $0$ is a parabolic fixed point (multiplier $1$) and $(0,\pi]$ lies in its
attracting direction.
\item There is a real-analytic $\Phi$ on $(0,\pi]$, unique up to an additive constant, with $\Phi\circ F=\Phi+1$ (the Fatou coordinate):
\[
\Phi(\theta)=\lim_{N\to\infty}\big[\Phi_0(F^N\theta)-N\big],\qquad \Phi_0(\theta)=\frac{3\pi}{\theta}+\frac32\log\theta,\qquad \Phi=\Phi_0+O(\theta).
\]
The coefficient $\tfrac32=1-b/a^2$ for $F=\theta-a\theta^2+b\theta^3+\dots$, $a=1/3\pi$, $b=-1/18\pi^2$, is the iterative residue
(r\'esidu it\'eratif), the formal conjugacy invariant of the germ \cite{Milnor,Ecalle}.
\item For every $\theta_0\in(0,\pi]$, $\theta_l=\Phi^{-1}(\Phi(\theta_0)+l)$ and
\[
\theta_l=\frac{3\pi}{\,l+\tfrac32\log l+\Phi(\theta_0)-\tfrac32\log(3\pi)+o(1)\,}.
\]
From orthogonal inputs ($\theta_0=\pi/2$) the constant is $\Phi(\pi/2)-\tfrac32\log3\pi=7.7548-3.3650=4.390$.
\end{enumerate}
\end{theorem}
\begin{proof}
(1) $1-J/\pi=\big(\pi(1-\cos\theta)-(\sin\theta-\theta\cos\theta)\big)/\pi=\theta^2u(\theta)$ with $u$ analytic, $u(0)=\tfrac12$, and
$F=2\arcsin\sqrt{(1-J/\pi)/2}=2\arcsin(\theta\sqrt{u/2})$, which is analytic. Expanding, $\cos F=1-\tfrac{\theta^2}2+\tfrac{\theta^3}{3\pi}+\tfrac{\theta^4}{24}
+O(\theta^5)$; matching $\cos(\theta+a_2\theta^2+a_3\theta^3)$ gives $a_2=-1/3\pi$ and $a_3=-a_2^2/2=-1/18\pi^2$. $F<\theta$ because
$J(\theta)/\pi-\cos\theta=(\sin\theta-\theta\cos\theta)/\pi>0$ on $(0,\pi]$.
(2) With $\Phi_0=1/(a\theta)+c\log\theta$ one computes $\Phi_0(F\theta)-\Phi_0(\theta)-1=\big((a^2-b)/a-ca\big)\theta+O(\theta^2)$. The choice
$c=(a^2-b)/a^2=\tfrac32$ makes the defect $O(\theta^2)$. Since $F^N\theta=O(1/N)$ the defects along an orbit are summable, so the limit exists,
is analytic (uniform limits of analytic functions on compact subsets of the petal), satisfies $\Phi\circ F=\Phi+1$, and differs from
$\Phi_0$ by $O(\theta)$. Uniqueness up to constants is the standard uniqueness of Fatou coordinates on a petal \cite[\S10]{Milnor}.
(3) Invert $\Phi_0(\theta_l)+O(\theta_l)=\Phi(\theta_0)+l$ with $\log\theta_l=\log3\pi-\log l+o(1)$.
\end{proof}
Numerically (Verification, item 3) $\Phi_0(\theta_l)-l$ equals $7.388$, $7.693$, $7.748$, $7.7542$, $7.7548$ at
$l=10,10^2,10^3,10^4,10^5$, while $3\pi/\theta_l-l$ keeps growing like $\tfrac32\log l$. At the competition depth the two laws give
\[
\theta_{16}=0.3779\ \text{(exact)},\qquad \frac{3\pi}{16+\frac32\log16+4.39}=0.3839,\qquad \frac{3\pi}{16}=0.5890 .
\]
So the stage-10 folding law is a leading-order statement whose correction is large at $L=16$.

\begin{corollary}[Gate counts, corrected]\label{cor:gates}
For two independent standard inputs ($\theta_0=\pi/2$ at large $n$) the expected number of layer-$l$ gates on which they differ is
$n\theta_{l-1}/\pi$. The total over $L$ layers is $\frac n\pi\sum_{k<L}\Phi^{-1}(\Phi(\pi/2)+k)=3n\log L+O(n)$, but the $O(n)$ term is large
and negative: for $L=16$ the total is $3.66\,n$, not $3n\log16=8.32\,n$ (the stage-10 prediction, which keeps only the leading term).
\end{corollary}

\paragraph{Farey, not Gauss.} The additive Rauzy--Veech induction for two intervals is the Farey map, which has a parabolic fixed point
with infinite invariant measure. Zorich's acceleration (grouping consecutive steps of one type) makes it the Gauss map, with a finite
measure and an integrable cocycle. The He-critical network is on the Farey side: depth acts on chords by a parabolic germ, and the
acceleration that turns it into a translation is precisely the Fatou coordinate. The network is not hyperbolic at any finite depth.
This is the renormalization-theoretic meaning of ``critical initialisation''.

\begin{theorem}[The two tiers are one flow]\label{thm:collapse}
Along Eldan's localization ($\Sigma_t=I/(1+t)$) two copies conditionally independent given $\mathcal F_t$ have input correlation
$t/(1+t)$; at infinite width the fraction of layer-$l$ gates on which they disagree is $\theta_{l-1}(t)/\pi$ with
$\theta_0(t)=\arccos\frac t{1+t}$, and $\E\Var(g_{l,a}\mid\mathcal F_t)=\theta_{l-1}(t)/2\pi$ (stage 10, Theorem \emph{depthtrickle}). Let
\[
\tau(l,t)=\Phi(\theta_0(t))+l-1\qquad\text{(the Fatou time).}
\]
\begin{enumerate}[leftmargin=*]
\item $\E\Var(g_{l,a}\mid\mathcal F_t)=\Phi^{-1}(\tau)/2\pi$: the unresolved gate variance depends on $(l,t)$ only through $\tau$. More generally
every infinite-width quantity built from finitely many exchangeable conditionally independent replicas at layer $l$ and time $t$ depends
on $(l,t)$ only through $\tau$.
\item The map $(t,l)\mapsto(t',l-1)$ with $\Phi(\theta_0(t'))=\Phi(\theta_0(t))+1$ preserves all of them: running the upper tier from $t$ to $t'$ is
the same as deleting one layer of the lower tier.
\item As $t\to\infty$, $\tau=3\pi\sqrt{(1+t)/2}-\tfrac34\log\frac{1+t}2+l-1+O(1)$, so the localization time worth one layer satisfies
$\Delta\sqrt{1+t}\to\sqrt2/3\pi=0.1500$.
\end{enumerate}
\end{theorem}
\begin{proof}
(1) $\theta_{l-1}(t)=F^{l-1}(\theta_0(t))=\Phi^{-1}(\Phi(\theta_0(t))+l-1)$ by Theorem~\ref{thm:fatou}. At infinite width the pre-activations of the
replicas at layer $l$ are jointly Gaussian with common variance (independent of $t$, since $X\sim\gamma$ marginally) and common pairwise
correlation $\cos\theta_{l-1}(t)$ (the dual-kernel recursion applied to each pair). (2) follows from (1). (3) $1-\cos\theta_0=1/(1+t)$
gives $\theta_0=\sqrt{2/(1+t)}\,(1+O(1/t))$; insert into $\Phi=\Phi_0+O(\theta)$.
\end{proof}
\paragraph{Putnam--Trevi\~no's Teichm\"uller flow, realized.} In \cite[\S12.8]{PT} the Teichm\"uller flow on bi-infinite diagrams built
from zippered rectangles is the identification
\[
(B,e^{-t_R}\nu_r,e^{t_R}\nu_s)\sim(\sigma B,\sigma_*\nu_r,\sigma_*\nu_s).
\]
It deforms the
state for a renormalization time $t_R$ and then shifts the diagram by one level. Theorem~\ref{thm:collapse}(2) is this identification
for the two-tier network. The diagram is the depth tower, its state is the gate law conditional on $\mathcal F_t$ (the upper tier), the shift
removes a layer, and the renormalization time is read off the Fatou coordinate. In log-time a layer costs
$\Delta\log(1+t)\approx2\sqrt2/(3\pi\sqrt{1+t})$: it becomes cheaper as the localization proceeds.

\paragraph{Data collapse.} The stage-10 measured disagreement fractions (width 512, 64 pairs, three times, six layers), sorted by
$\tau$ (Verification, item 4):
\begin{center}\small
\begin{tabular}{rrrcc|rrrcc}\toprule
$\tau$ & $t$ & $l$ & meas. & pred. & $\tau$ & $t$ & $l$ & meas. & pred.\\\midrule
9.80&1&1&.333&.333 & 24.80&1&16&\textbf{.136}&.115\\
10.80&1&2&.287&.292 & 27.02&9&8&.106&.105\\
12.80&1&4&.223&.236 & 31.02&9&12&.082&.092\\
16.80&1&8&.174&.173 & 35.02&9&16&\textbf{.098}&.081\\
20.02&9&1&.140&.144 & 63.76&99&1&.044&.045\\
20.80&1&12&.130&.138 & 70.76&99&8&.040&.041\\
21.02&9&2&.135&.136 & 74.76&99&12&.033&.039\\
23.02&9&4&.116&.124 & 78.76&99&16&\textbf{.043}&.037\\\bottomrule
\end{tabular}
\end{center}
Points from different localization times interleave on a single decreasing curve, e.g.\ $(t,l)=(9,1)$ at $\tau=20.0$ and $(1,12)$ at $\tau=20.8$.
The three last-layer points (bold) all sit above it. This is the finite-width drift of the last layer that stage 10 recorded, and the
collapse makes it a clean diagnostic: a mean-field (replica-symmetric, infinite-width) quantity cannot leave the $\tau$-curve, so the
excess at $l=16$ is a finite-width effect.

\begin{proposition}[Tail equivalence at infinite width]\label{prop:tail}
For $\theta_0>0$, $\sum_l\theta_l=\infty$ (harmonic, from Theorem~\ref{thm:fatou}). At infinite width two distinct inputs therefore disagree on a
positive fraction $\theta_{l-1}/\pi$ of the gates of every layer, and the total expected disagreement diverges even per neuron. The
gate sequences of distinct inputs are never tail-equivalent.
\end{proposition}
\begin{conjecture}[Finite width is a hyperbolic perturbation]\label{conj:hyper}
At width $n$ the projective contraction of products of gated Gaussian matrices adds a linear term,
$\theta\mapsto(1-c_n)\theta-\theta^2/3\pi+\dots$ with $c_n\asymp1/n$ (the gap between the top two Lyapunov exponents). The parabolic point becomes
attracting-hyperbolic, the crossover depth is $\asymp n$, $\sum_l\theta_l<\infty$, and almost every pair of inputs (on non-dead
networks) becomes eventually tail-equivalent. This would settle the stage-10 question about the tail algebra of the depth tower: it is
decided by the summability of the Fatou orbit. The competition networks ($L=16\ll n=1024$) sit deep in the parabolic regime.
\end{conjecture}

\section{The walls as a stratified singular foliation}
\label{sec:walls}
Work on $\R^n\setminus\{0\}$; everything is conical. The depth-$L$ tiling stratifies input space into relatively open polyhedral cones:
open tiles, open pieces of walls, and so on. These are the \emph{faces}. Each face is convex, hence contractible.
\begin{definition}[Wall foliation]
The wall foliation $\cF_L$ is the module of (formal) vector fields tangent to every wall $Z_{l,a}$, $l\le L$. Its leaves are the faces.
At a point of a face $S$ of codimension $j$ the \emph{transverse model} is the induced formal foliation of a $j$-dimensional slice.
\end{definition}
The radial (Euler) field $x\cdot\nabla$ is tangent to every wall, because walls are cones. So the drift of the Ornstein--Uhlenbeck generator
$\Delta-x\cdot\nabla$ is longitudinal and only the diffusion is transverse. The Euler--Stein identity $\E f=\E[x\cdot\nabla f]=\E[\Delta f]$ for
$1$-homogeneous $f$ records where the longitudinal part meets the transverse one.

\begin{theorem}[Normal and bent crossings]\label{thm:bent}
Let $S$ be a codimension-two face at a generic point of which exactly two walls meet, $Z=Z_{l,a}$ and $Z'=Z_{k,b}$ with $k\le l$. Let
$T=\partial z_{l,a}/\partial h_{k,b}$ be the gated transport coefficient (intermediate gates fixed at the point; $T=0$ if $k=l$).
\begin{enumerate}[leftmargin=*]
\item If $T=0$ the walls are transverse hyperplanes near $S$. The transverse model is the log foliation of $y_1y_2=0$, generated by $y_1\partial_1$
and $y_2\partial_2$; the isotropy Lie algebra is $\R^2$ (abelian). The foliation is Debord near $S$ (its algebroid is the log tangent
bundle), its Nash blowup is trivial, and its Helffer--Nourrigat cone is the whole log cotangent bundle \cite{Louis}. The holonomy groupoid
is the $b$-groupoid of Melrose's calculus \cite{Melrose}.
\item If $T\ne0$, choose transverse coordinates $y_1=z_{k,b}$ and $y_2=$ the linear function equal to $z_{l,a}$ on the side $y_1<0$. Then
$z_{l,a}=y_2+T\relu(y_1)$ near $S$, so $Z=\{y_2=0,\ y_1\le0\}\cup\{y_2+Ty_1=0,\ y_1\ge0\}$ is \emph{bent} along $Z'$. The formal transverse model is
the log foliation of the central three-line arrangement
\[
y_1\,y_2\,(y_2+Ty_1)=0 ,
\]
a free module (Saito) with basis $E=y_1\partial_1+y_2\partial_2$ and $\vartheta_T=-y_1(y_2+Ty_1)\partial_1$. The isotropy Lie algebra is
$\mathrm{span}(E,\vartheta_T)$ with $[E,\vartheta_T]=\vartheta_T$, i.e.\ $\aff(1)$, non-abelian, and the linear isotropy is $\R E$, of dimension $1$
instead of $2$.
\item The integrated isotropy at a bent crossing is the connected $ax+b$ group $\mathrm{Aff}_+(1)$, which Francis identifies with the jet
group $J^2$ \cite{Francis}. The groupoid of $\cF_L$ restricted to $S$ is $S\times S\times\mathrm{Aff}_+(1)$, with $C^*$-algebra $C^*(\mathrm{Aff}_+(1))\otimes\mathcal K$, an
extension of $C_0(\R)$ by $\mathcal K\oplus\mathcal K$. At a normal crossing it is $C_0(\R^2)\otimes\mathcal K$.
\end{enumerate}
\end{theorem}
\begin{proof}
Near a generic point of $S$ only the gate $g_{k,b}$ changes, and crossing $Z'$ toggles $h_{k,b}$ between $0$ and $z_{k,b}$, which changes
$z_{l,a}$ by $T\,z_{k,b}$. This gives the normal form. A formal vector field is tangent to a half-plane through $S$ iff each homogeneous
component is, and a homogeneous polynomial field tangent to a half-line in the slice is tangent to the whole line (polynomial identity on
a ray). So the formal module is the module of fields tangent to the three lines $\{y_1=0\}$, $\{y_2=0\}$, $\{y_2+Ty_1=0\}$; the extensions of the
two half-lines of $Z$ are \emph{phantom} lines that the formal structure cannot distinguish from the wall. Saito's criterion holds:
$\det(E,\vartheta_T)=y_1y_2(y_2+Ty_1)$. $[E,\vartheta_T]=(2-1)\vartheta_T$ because $\vartheta_T$ is homogeneous of degree $2$, and $\vartheta_T\notin I_0\cF$ (otherwise
$\vartheta_T=\ell E$ with $\ell$ linear and the determinant would vanish). A linear field tangent to three distinct lines through $0$ has three
eigenlines, hence is scalar. For $T=0$ two of the lines coincide and the reduced arrangement is $y_1y_2=0$. (3): the free module is the
section module of a Lie algebroid with injective anchor on the regular part, so the holonomy groupoid is the Lie groupoid integrating it
\cite{AS,Louis}; the simply connected group with Lie algebra $\aff(1)$ is $\mathrm{Aff}_+(1)$, which has trivial centre. The structure of
$C^*(\mathrm{Aff}_+(1))$ is classical.
\end{proof}
The check (Verification, item 5) solves the tangency conditions for homogeneous polynomial fields in the slice. The dimensions in degrees
$1,2,3$ are $2,4,6$ at $T=0$ and $1,3,5$ at $T=0.3$, and $[E,\vartheta_T]=\vartheta_T$ holds symbolically.

\paragraph{The bent crossings are the memory.} $T$ is the coefficient with which the kink of neuron $(k,b)$ enters the deep pre-activation:
the coupling at which Price's theorem gives birth to the third cumulant (stage 10, Theorem \emph{birth}). In the simplest case it is literally
the cumulant. For $z_{2,a}=\sum_bW_2(a,b)\relu(w_b\cdot x)$ with orthogonal first-layer rows, the $w_b\cdot x$ are independent, the bending
coefficients are $T_{ab}=W_2(a,b)$, and
\[
\kappa_3(z_{2,a})=c_3\sum_bT_{ab}^3|w_b|^3,\qquad c_3=\kappa_3(\relu Z)=\frac{\pi+2}{(2\pi)^{3/2}}=0.3265 .
\]
So $z_{2,a}$ is Gaussian iff its wall is nowhere bent, and its skewness is the sum of cubed bendings over its bent crossings. With
correlated rows, cross terms appear, weighted by the codegrees (cosines) of the normal crossings among the shallow walls. These are
the walk sums with exact contacts of stage 10, now read as an expansion over the codimension-two strata of the wall foliation: bent
crossings are vertices with a contact, normal crossings are edges. The Gaussian closure is exact for any wall system with only normal
crossings, and every family of jointly Gaussian pre-activations is the pre-activation family of such a system in an augmented Gaussian
space (realize the covariance as $\tilde z_l=M_l\xi$). Its error is therefore carried by the bent crossings.

\paragraph{Smooth versus piecewise linear.} The piecewise-linear homeomorphism $(y_1,y_2)\mapsto(y_1,y_2+T\relu(y_1))$ unbends $Z$ into a
normal crossing. In the PL category every crossing is normal. The bending, the phantom line and the non-abelian isotropy exist only for
the smooth (formal) structure, and the Gaussian measure is a smooth object, not a PL one. This is the geometric content of ``Gaussian
calculus sees the kinks'': non-Gaussianity is the obstruction to smoothing the unbending chart.

\begin{proposition}[No holonomy]\label{prop:noholonomy}
(1) Every face is contractible, so by Fischer--Laurent-Gengoux (\cite{FLG}, Cor.~3.10) the formal wall foliation along any face is the
product of the face with its transverse model: outer holonomy and principal-bundle data are trivial. (2) Every wall $Z_{l,a}$ has the
global defining function $z_{l,a}$; for each $k$ the transverse-order-$k$ structure it induces is Scott's $b^k$ structure with a global
$(k-1)$-jet \cite{Scott}, whose holonomy invariant is trivial (\cite{Francis}, \S1). Hence the noncommutativity of the wall foliation is purely
isotropic: abelian at normal crossings, $\aff(1)$ at bent crossings, and solvable at higher-codimension faces.
\end{proposition}
The contrast with tilings is instructive. In a substitution tiling the boundary relations $R^j_B$ of Julien--Savinien are large: dense and
of tail type. In the network every face is a convex cone and the global defining functions kill all holonomy. Whatever the network has
that is genuinely noncommutative lives at points (isotropy) and in the depth-and-localization dynamics (Sections~\ref{sec:parabolic},
\ref{sec:deformation}), not in the leaf space.

\section{The upper tier as a semiclassical deformation}
\label{sec:deformation}
The upper tier is the family of Gaussian laws $N(m,\hbar I)$, $\hbar\in(0,1]$; Eldan's path has $\hbar=1/(1+t)$ and a moving centre. For a
network output $f$ put
\[
F(m,\hbar)=\E f(m+\sqrt\hbar Z),\qquad \partial_\hbar F=\tfrac12\Delta_mF,\qquad u=F(0,1).
\]
In Paul's language $N(m,\hbar I)$ is the squared modulus of a Gaussian coherent state of width $\sqrt\hbar$, and $F(\cdot,\hbar)$ is the Husimi
(anti-Wick, Toeplitz lower) symbol of the operator of multiplication by $f$. The classical limit $\hbar\to0$ returns $f$, which is
singular along the wall foliation.

\begin{theorem}[Boundary layer; the parabolic blowup]\label{thm:layer}
Let $m$ lie within distance $\rho$ of a single flat wall piece $W$ of $f$, with all other walls at distance at least $R>\rho$, signed distance
$d$ to $W$ and kink $\kappa$ (the jump of the normal derivative of $f$ across $W$). Then
\[
F(m,\hbar)=f(m)+\sqrt\hbar\,\kappa\,\psi\!\Big(\frac d{\sqrt\hbar}\Big)+O\big(e^{-(R-\rho)^2/2\hbar}\big),
\]
\[
\psi(u)=\varphi(u)-|u|\,\Phi_{\mathcal N}(-|u|)=\E(u+Z)^+-u^+ .
\]
Hence $F$ lifts to a smooth function on the parabolic blowup of $W\times\{0\}$ in $(m,\hbar)$-space, with radial coordinate
$r=(d^2+\hbar)^{1/2}$ and angular coordinate $d/\sqrt\hbar$. The correction is homogeneous of degree one in $r$, i.e.\ of transverse
order one in Francis's sense. Near a codimension-$j$ face, iterated parabolic blowups are needed, performed deepest stratum first:
the combinatorics of the resolution is the face multi-diagram of Section~\ref{sec:borders}.
\end{theorem}
\begin{proof}
On the ball of radius $R$ about the foot point, $f=\ell+\kappa\,d^+$ with $\ell$ linear. $\E\ell(m+\sqrt\hbar Z)=\ell(m)$ and
$\E(d+\sqrt\hbar Z_\nu)^+=\sqrt\hbar\,(\varphi(u)+u\Phi_{\mathcal N}(u))$ with $u=d/\sqrt\hbar$; subtract $d^+$. The contribution of the complement of the ball
is bounded by the Gaussian tail times a linear bound on $|f|$.
\end{proof}
The check (Verification, item 6) uses a two-dimensional network with six first-layer walls and a base point $3.9$ from the other walls.
For $\hbar\in\{0.01,0.05,0.2,0.5\}$ and $d/\sqrt\hbar\in\{0,0.7\}$ the formula holds to relative error $10^{-8}$, the quadrature accuracy.
Stage 10's Euler--Stein constant reads naturally on this blowup. A defect of Hermite contact order $d$ is homogeneous of degree $d+1$ in
$r$ (transverse order $d+1$), and averaging over a uniformly distributed radius gives $\int_0^1r^{d+1}dr=1/(d+2)$.

\begin{proposition}[The target is deep in the quantum regime; the right limit is kinetic]\label{prop:kinetic}
(1) At scale $\hbar$ the expected number of unresolved first-layer gates is $\E\sum_a\Var(g_{1,a}\mid\text{ball})=n\,\theta_0(\hbar)/2\pi$ with
$\theta_0=\arccos(1-\hbar)$. At the target $\hbar=1$ this is $n/4$, and the single-wall (semiclassical) regime $n\theta_0/2\pi\lesssim1$ requires
$\hbar\lesssim2\pi^2/n^2\approx1.9\times10^{-5}$ for $n=1024$. (2) Each wall bends the downstream pre-activations by $T=O(n^{-1/2})$ and there are
$n$ walls per layer with $\sum_bT_{ab}^2g_b\approx1$ at He initialisation: coupling$^2\times$density$=O(1)$. This is the weak-coupling
(van Hove) scaling of a particle among many weak scatterers \cite{EY}. The infinite-width (NNGP, moment-chain) description is the
kinetic limit of this scaling, not a semiclassical expansion, and the finite-width memory is its correction. Ehrenfest-type arguments
therefore cannot control the memory.
\end{proposition}
\begin{proof}
(1) is stage 10's unresolved-variance formula at layer one with correlation $1-\hbar$; at $\hbar=1$ each gate is a fair coin. For small
$\hbar$, $\theta_0\approx\sqrt{2\hbar}$. (2) $W_2(a,b)\sim N(0,2/n)$ and half the gates are open.
\end{proof}

\paragraph{Paul's two-point symbols.} Past the Ehrenfest time Paul's symbols become kernels $\sigma(z,z')$ along unstable leaves, and the
matching states are off-diagonal Husimi functions $R_\psi(z,z')$. The network's counterparts are stage 10's time-correlated Kikuchi weights
$K_{l,t}(a,b)=\mathbb P(g_{l,a}(X)=1,g_{l,b}(X')=1)$ and the angle operator $\Theta_{l,t}=E_{\mathfrak B_l}E_{\mathfrak F_t}E_{\mathfrak B_l}$. Both are built
from two coherent samples $X,X'$ of one localized law, so they are two-point objects of exactly Paul's kind. In the network the role of the
unstable foliation (the direction in which the coherent state is sheared) is played by the folding of the depth tower, whose clock is
the Fatou time. Theorem~\ref{thm:collapse} says these two-point objects are functions of $\tau$ alone at infinite width. (Analogy, with
the exact content of Theorem~\ref{thm:collapse}.)

\paragraph{Giannakis--Montgomery's Fock space.}
\begin{enumerate}[leftmargin=*]
\item The Wiener chaos of $\gamma$ is the symmetric Fock space over $\R^n$. Exponential vectors are the tilted densities $e^{m\cdot x-|m|^2/2}$
(coherent states), with $a_ve^{m\cdot x-|m|^2/2}=\langle v,m\rangle e^{m\cdot x-|m|^2/2}$: localization centres are coherent states. A moment
chain of order $K$ is a truncation to $\le K$ quanta. The linear part of a layer acts by second quantization ($k$-th moment tensors
transform by $W^{\otimes k}$, the functor $\Gamma(W)$), and the \textsc{ReLU} is a vertex that mixes grades.
\item Depth drives the network toward the commutative quotient. Stage 10's Mehler pair temperature gives a mean token number
$\approx1/(1-\cos\theta_{l-1})$ for a neuron pair at layer $l$, so the effective Planck constant is $\hbar_l\approx1-\cos\theta_{l-1}\approx\theta_{l-1}^2/2
\approx9\pi^2/2(l+\tfrac32\log l+4.39)^2$ by Theorem~\ref{thm:fatou}. On high Fock grades the weighted canonical commutation relations become
commutative (\cite{GM}, Prop.~32), which is the classical limit.
\item The target is a kernel mean embedding: $u=\E_\gamma h_L$ is the mean embedding of $\gamma$ for the random feature map $h_L$, with kernel
$k_L(x,x')=h_L(x)\cdot h_L(x')/n$ (exact). For any surrogate law $\nu$, $\frac1n\|\E_\gamma h_L-\E_\nu h_L\|^2=\mathrm{MMD}^2_{k_L}(\gamma,\nu)$
\cite{MMD}. The competition's mean squared error of such an estimator is a maximum mean discrepancy in the depth-$L$ feature space.
\item The Fibonacci dilation realizes a non-normal generator as the corner $(p_1\otimes1)L^kE_n$ of a skew-adjoint banded operator on
$\ell^2(\mathbb N)\otimes H$. Stage 10 realized the network as the corner $[(1-D_\downarrow P_x)^{-1}]_{L0}$ of a resolvent of a banded chiral
operator on $\ell^2(\{0,\dots,L\})\otimes\R^n$. A unitary dilation of the gated transfer, which would make the layer law's Gaussian part
invariant (a Koopman--von Neumann wavefunction for the network), is open.
\end{enumerate}

\section{Borders: the mean lives on the codimension-one faces}
\label{sec:borders}
\begin{proposition}[Two Bratteli diagrams of the depth tower]\label{prop:twodiagrams}
\begin{enumerate}[leftmargin=*]
\item \emph{Tile diagram.} The vertices at level $l$ are the nonempty depth-$l$ tiles, edges are inclusions, and $\nu(v)=\gamma(v)$ is a one-sided
state ($\nu(v)=\sum_{\text{children}}\nu$). It is a tree, so its AF algebra is commutative. On a great circle, with the angular order as
$\le_s$, it is Putnam--Trevi\~no's ``decimal expansion'' picture: the circle is the quotient of the path space by successor
identifications, and the successor pairs are the two sides of the breakpoints.
\item \emph{Fiber diagram.} The supertiles of level $l$ are the fibres of $h_l$; write $x\approx_lx'$ iff $h_l(x)=h_l(x')$. These relations
increase with $l$ (the AF direction), their classes are polyhedra of dimension $n-\operatorname{rank}(D_lW_l\cdots D_1W_1)$ (about $n/2$ from the
first layer on), and $\approx_\infty=\bigcup_l\approx_l$ (eventual indistinguishability) is contained in the tail relation of the \emph{letter
diagram}, whose level-$l$ vertices are the layer-$l$ gate patterns alone (finite rank $2^n$). The letter diagram's tail relation is the network's
$R_{AF}$.
\item \emph{Borders.} Call $n-j$ the border dimension of $x$ if, for every large depth, $x$ lies on a face of codimension $j$. At infinite width
and infinite depth every Gaussian ball, of any radius $\sqrt\hbar>0$, meets walls of every layer, with $\sum_ln\,\theta_{l-1}(\hbar)/2\pi=\infty$ by
Theorem~\ref{thm:fatou}. So points of border dimension $<n$ are dense; they are null. This is Julien--Savinien's thin set (their
Prop.~3.22 and Rem.~3.26), and it is dense precisely because depth is parabolic: under the hyperbolic finite-width dynamics of
Conjecture~\ref{conj:hyper} the sum converges.
\item \emph{Births.} A wall of layer $l$ inside a tile of depth $l-1$ is a face that does not descend from a face of the parent. The face
multi-diagram (Bratteli diagrams of faces of each codimension, joined by incidence edges) therefore has \emph{birth} edges from tiles at
level $l-1$ to codimension-one faces at level $l$. These are the reverse of Julien--Savinien's escaping edges (faces swallowed by
supertiles), and they are where stage 10 locates the creation of contacts.
\end{enumerate}
\end{proposition}
\begin{proof}
(1) and (2) are immediate from the definitions; $h_l(x)=h_l(x')$ forces equal gates at all later layers. (3) uses the infinite-width
unresolved variance at every layer and the divergence of $\sum\theta_l$. (4) is the definition of a wall.
\end{proof}

\begin{proposition}[The mean is a border functional]\label{prop:border}
For any network output $f$ (piecewise linear, $1$-homogeneous),
\[
\E_\gamma f=\E\Delta f=\frac1{\sqrt{2\pi}}\sum_{F}\kappa_F\,\alpha_F ,
\]
where $F$ runs over the codimension-one faces of the depth-$L$ tiling, $\kappa_F$ is the jump of the normal derivative across $F$, and $\alpha_F$
is the Gaussian measure of the cone $F$ inside its hyperplane (its solid-angle fraction). Equivalently, let $H_\partial=L^2(\bigcup_FF\times\{+,-\},\varphi\,dS)$
(two copies of each wall piece, one per side), let $\Gamma$ be the side grading, $\mathbb F$ the side flip, and $\pi(g)$ multiplication by the
one-sided values of a field $g$ that is constant on open tiles. Then
\[
[\mathbb F,\pi(g)]=0\iff g\ \text{is continuous across every wall},\qquad
\E_\gamma f=\tfrac12\Tr_\varphi\big(\Gamma\,\mathbb F\,[\mathbb F,\pi(\partial_\nu f)]\big).
\]
\end{proposition}
\begin{proof}
Euler and Stein give $\E f=\E[x\cdot\nabla f]=\E\Delta f$, and $\Delta f=\sum_F\kappa_F\,\delta_F$ as a distribution. Moreover $\int_F\varphi_n\,dS=\varphi_1(0)\,\alpha_F$
because $F$ is a cone through $0$. On one face, with $\pi(g)=\diag(g_+,g_-)$, $\Gamma\mathbb F[\mathbb F,\pi(g)]=\diag(g_+-g_-,g_+-g_-)$; take $g=\partial_\nu f$.
\end{proof}
The second form is Putnam--Trevi\~no's Fredholm module with the successor flip replaced by the wall flip. Their Theorem 9.7 (functions
descend to the quotient iff they commute with the flip) becomes continuity across walls, and the Gaussian mean of the network becomes the
supertrace pairing of the flip with the normal derivative: the kink current of stage 10. On a continuum of wall points the commutators
are multiplication operators rather than compact ones, so this is a Gaussian-weighted analogue, not a Fredholm module in the strict
sense. Its restriction to a great circle, which has finitely many breakpoints, is one. The two-dimensional check (Verification,
item 7) has $70$ breakpoints and gives $\E f=0.36919251$ both by angular quadrature and by the border sum. In Julien--Savinien's terms, the
whole mean is carried by the border paths of the codimension-one face diagram: the generic (AF) points contribute nothing.

\section{What this explains, predicts, and leaves open}
\label{sec:outlook}
\subsection{What the placement explains}
\begin{enumerate}[leftmargin=*]
\item \textbf{Why the folding law needed a correction.} Depth is a parabolic germ with iterative residue $3/2$, so $3\pi/l$ misses a
$\tfrac32\log l$ and a constant that are large at $l=16$ (Theorem~\ref{thm:fatou}). The He-critical network is Farey-like, not
Gauss-like.
\item \textbf{Why localization resolves deep gates first, quantitatively.} Depth and localization time are two coordinates of one
renormalization flow; one layer is worth $\sqrt2/3\pi$ of $\sqrt{1+t}$ (Theorem~\ref{thm:collapse}). This is the Teichm\"uller flow of
Putnam--Trevi\~no: a state deformation followed by a shift of the diagram.
\item \textbf{Where the memory lives geometrically.} At bent crossings of the wall foliation, where the isotropy is $\aff(1)$ instead of
abelian and the bending is the gated transport coefficient (Theorem~\ref{thm:bent}). The Gaussian closure is exact for wall systems with
only normal crossings.
\item \textbf{Why tiling-theoretic holonomy is absent.} Faces are convex and walls have global defining functions
(Proposition~\ref{prop:noholonomy}). The noncommutativity that tiling theory finds in leaf spaces is, for the network, in isotropy and in
the two-tier dynamics.
\item \textbf{Why semiclassical expansions cannot reach the target.} At $\hbar=1$ there are $n/4$ unresolved gates per layer; the
infinite-width description is a kinetic (van Hove) limit (Proposition~\ref{prop:kinetic}).
\item \textbf{Why the mean is a boundary quantity.} It is the border functional of the codimension-one face diagram, or the supertrace
pairing of the wall flip with the normal derivative (Proposition~\ref{prop:border}).
\end{enumerate}

\subsection{Predictions (falsifiable)}
\begin{enumerate}[leftmargin=*]
\item Two independent standard inputs differ on $n\theta_{l-1}/\pi$ gates of layer $l$, where $\theta_1=1.247$, $\theta_2=1.054$, and so on down to
$\theta_{16}=0.378$; the total over $16$ layers is $3.66\,n$. This replaces stage 10's $3n\log L=8.32\,n$ (Corollary~\ref{cor:gates}).
\item Every infinite-width two-tier quantity collapses on the Fatou time $\tau$. The departure of finite-width measurements from the
$\tau$-curve is a finite-width (memory) signal; it is visible at the last layer for all three localization times in the stage-10 data.
\item With orthogonal first-layer rows, $\kappa_3(z_{2,a})=c_3\sum_bW_2(a,b)^3|w_b|^3$ exactly, and $z_{2,a}$ is Gaussian iff no crossing is bent.
\item (Conjectural) At width $n$ the parabolic regime ends at depth $\asymp n$. Beyond it the angle decays exponentially, the walls met by a
small ball become finite in number, and gate sequences become eventually tail-equivalent (Conjecture~\ref{conj:hyper}).
\end{enumerate}

\subsection{Open problems}
\begin{enumerate}[leftmargin=*]
\item \textbf{The finite-width linear term.} Compute the projective Lyapunov gap of products of gated He matrices, prove
Conjecture~\ref{conj:hyper}, and decide the tail algebra of the letter diagram (the stage-10 question).
\item \textbf{The error of the Gaussian closure as a sum over bent crossings.} Interpolate each bent crossing to a normal one (send its
$T$ to $0$) and write the closure error as a sum of crossing-local terms. This is the foliation form of stage 10's contact-resolved
estimator problem, and the $\aff(1)$ isotropy suggests that each term is a commutator.
\item \textbf{$K$-theory and index.} The wall-foliation algebra has a finite composition series over the faces with subquotients
$C^*(\text{isotropy})\otimes\mathcal K$ ($C_0(\R^j)$ at normal crossings, $C^*(\mathrm{Aff}_+(1))$ at bent ones). Is there an index pairing, sharper than
the supertrace formula, whose value is the mean or one of its chain corrections? Savinien--Bellissard's spectral sequence for tilings is
the template.
\item \textbf{Nonnegative networks as flat surfaces.} For $W\ge0$ (e.g.\ the majorant $|W|$) Putnam--Trevi\~no's construction yields a
genuine translation surface for each input, with widths $h$, heights $\delta$ and area $f$. Do its Teichm\"uller flow and asymptotic cycle bound
the signed network?
\item \textbf{A Koopman--von Neumann wavefunction.} A unitary dilation of the gated transfer in the spirit of Giannakis--Montgomery,
under which the Gaussian part of the layer law is invariant.
\item \textbf{Chains in Fatou time.} Stage 9 and stage 10 transport Euler--Stein constants and window corrections layer by layer. If their
depth dependence is a function of $\tau$ rather than of $l$, one calibration at one depth fixes all depths.
\end{enumerate}

\section*{Verification}
The checks are run by the script \texttt{verify\_stage11.py};
its output is saved next to this note.
\begin{center}\small
\begin{tabular}{>{\raggedright\arraybackslash}p{5.4cm}>{\raggedright\arraybackslash}p{9.4cm}}\toprule
identity & result\\\midrule
1. $\langle h_l,\delta_l\rangle=f$ at every layer; rescaling & $n=64$, $L=8$: max error $2.9\times10^{-15}$; per-neuron rescaling changes $f$ by
$4\times10^{-16}$ and the rectangles by $9\times10^{-16}$\\
2. Rauzy--Veech as a gated linear map & $d=4$, $\pi=(4321)$, 12 steps: integer Jacobian, determinant $1$; $186/200$ perturbations stay in the
cylinder and obey the linear map to $2\times10^{-11}$; all $64$ type patterns of length $6$ occur\\
3. Parabolic germ, Fatou coordinate & $(F-\text{series})/\theta^4=-0.00766,-0.00767$ at $\theta=10^{-2},10^{-3}$; $\Phi_0(\theta_l)-l=7.388,7.693,7.748,7.7542,7.7548$
at $l=10,\dots,10^5$; the $3\pi/l$ law errs by $84\%,11\%,1.5\%$ at $l=10,10^2,10^3$\\
4. Two-tier collapse & stage-10 width-512 data monotone in $\tau$ except the three last-layer points; $\Delta\sqrt{1+t}$ per layer
$=0.1533,0.1511,0.1504$ at $t=99,999,9999$ (limit $\sqrt2/3\pi=0.1500$)\\
5. Bent crossing & tangent fields of degree $1,2,3$: $2,4,6$ at $T=0$, $1,3,5$ at $T=0.3$; $[E,\vartheta_T]=\vartheta_T$ symbolically\\
6. Boundary layer & relative error $2$--$8\times10^{-8}$ for $\hbar\in\{0.01,0.05,0.2,0.5\}$, $d/\sqrt\hbar\in\{0,0.7\}$\\
7. Mean as a border functional & $\R^2$ input, $16+16$ hidden units, $70$ breakpoints: $0.36919251$ by quadrature and by the border sum\\\bottomrule
\end{tabular}
\end{center}
The exact folding angles from $\theta_0=\pi/2$ are $\theta_1,\dots,\theta_{16}=1.247$, $1.054$, $0.921$, $0.822$, $0.744$, $0.681$, $0.628$, $0.584$, $0.546$, $0.513$, $0.483$,
$0.458$, $0.435$, $0.414$, $0.395$, $0.378$, and $\sum_{l=1}^{16}\theta_{l-1}/\pi=3.658$.

\begin{thebibliography}{99}\small
\bibitem{stage10} Stage 10 note: \emph{the selected Dirac operator} (this repository, \texttt{notes/stage10}).
\bibitem{stage9} Stage 9 note: \emph{the mean of a random ReLU network as a central localization} (\texttt{notes/stage9}).
\bibitem{BJS} J.~Bellissard, A.~Julien, J.~Savinien, Tiling groupoids and Bratteli diagrams, Ann.~Henri Poincar\'e 11 (2010) 69--99; arXiv:0911.0080.
\bibitem{JS} A.~Julien, J.~Savinien, Tiling groupoids and Bratteli diagrams II: structure of the orbit equivalence relation, arXiv:1005.2965.
\bibitem{PT} I.~F.~Putnam, R.~Trevi\~no, Bratteli diagrams, translation flows and their $C^*$-algebras, arXiv:2205.01537.
\bibitem{GM} D.~Giannakis, M.~Montgomery, Measure-free Koopman--von Neumann dynamics and noncommutative geometry, arXiv:2608.11591.
\bibitem{Paul} T.~Paul, Semiclassical approximation and noncommutative geometry, preprint hal-00617372 (2011).
\bibitem{Francis} M.~Francis, On singular foliations tangent to a given hypersurface, arXiv:2311.03940.
\bibitem{FLG} S.-R.~Fischer, C.~Laurent-Gengoux, A classification of neighborhoods around leaves of a singular foliation, arXiv:2401.05966.
\bibitem{Louis} R.~Louis, On longitudinal differential operators and Nash blowups, arXiv:2509.01133.
\bibitem{AS} I.~Androulidakis, G.~Skandalis, The holonomy groupoid of a singular foliation, J.~reine angew.~Math.~626 (2009) 1--37.
\bibitem{AMY} I.~Androulidakis, O.~Mohsen, R.~Yuncken, A pseudodifferential calculus for maximally hypoelliptic operators and the Helffer--Nourrigat conjecture, arXiv:2201.12060.
\bibitem{Saito} K.~Saito, Theory of logarithmic differential forms and logarithmic vector fields, J.~Fac.~Sci.~Univ.~Tokyo 27 (1980) 265--291.
\bibitem{Melrose} R.~B.~Melrose, \emph{The Atiyah--Patodi--Singer Index Theorem}, A~K~Peters 1993.
\bibitem{Scott} G.~Scott, The geometry of $b^k$ manifolds, J.~Symplectic Geom.~14 (2016) 71--95.
\bibitem{Veech} W.~A.~Veech, Gauss measures for transformations on the space of interval exchange maps, Ann.~of Math.~115 (1982) 201--242.
\bibitem{Zorich} A.~Zorich, Finite Gauss measure on the space of interval exchange transformations, Ann.~Inst.~Fourier 46 (1996) 325--370.
\bibitem{Milnor} J.~Milnor, \emph{Dynamics in One Complex Variable}, 3rd ed., Princeton 2006 (\S10, parabolic points and Fatou coordinates).
\bibitem{Ecalle} J.~\'Ecalle, \emph{Les fonctions r\'esurgentes}, Publ.~Math.~Orsay 1981; S.~M.~Voronin, Analytic classification of germs of conformal mappings $(\mathbb C,0)\to(\mathbb C,0)$, Funct.~Anal.~Appl.~15 (1981) 1--13.
\bibitem{CS} Y.~Cho, L.~K.~Saul, Kernel methods for deep learning, NeurIPS 2009.
\bibitem{EY} L.~Erd\H{o}s, H.-T.~Yau, Linear Boltzmann equation as the weak coupling limit of a random Schr\"odinger equation, Comm.~Pure Appl.~Math.~53 (2000) 667--735.
\bibitem{MMD} B.~K.~Sriperumbudur, A.~Gretton, K.~Fukumizu, B.~Sch\"olkopf, G.~Lanckriet, Hilbert space embeddings and metrics on probability measures, JMLR 11 (2010) 1517--1561.
\bibitem{Eldan} R.~Eldan, Thin shell implies spectral gap up to polylog via a stochastic localization scheme, GAFA 23 (2013) 532--569.
\end{thebibliography}
\end{document}
